% concatenated from List_distinguishing_chromatic_number_of_graphs_containing_only_small_complete_bigraphs.tex
\documentclass[12pt,reqno,a4paper]{amsart}
%%%%%%%%%%%%%%%%%%packages
\usepackage{hyperref}
\usepackage{enumerate}
\usepackage[left=0.75in, right=0.75in, top= 0.9in, bottom=1.65in]{geometry}
\usepackage{amsmath,amssymb,amsthm,amsfonts}
\setlength{\parindent}{0pt}
\setlength{\parskip}{0.45em}

 \linespread{1.15}

\usepackage{graphicx}
\usepackage{tikz}

%\extrapackages
\usetikzlibrary{decorations.pathreplacing}
\usetikzlibrary{decorations.markings}
\usetikzlibrary{calc} % Needed for ($(v0)!0.5!(v4)$) syntax


\usetikzlibrary{graphs, graphs.standard}

\tikzset{
	modal/.style={>=stealth’,shorten >=1pt,shorten <=1pt,auto,node distance=1.5cm,
		semithick},
	world/.style={circle, draw,minimum size=.1cm,fill=gray!15},
	point/.style={circle,draw,inner sep=0.3mm,fill=black},
	circ/.style={circle,draw,inner sep=0.1mm,fill=white},
	reflexive above/.style={->,loop,looseness=7,in=120,out=60},
	reflexive below/.style={->,loop,looseness=7,in=240,out=300},
	reflexive left/.style={->,loop,looseness=7,in=150,out=210},
	reflexive right/.style={->,loop,looseness=7,in=30,out=330}
}


\usetikzlibrary{shapes}
\usetikzlibrary{plotmarks}
\usetikzlibrary{arrows}
\usetikzlibrary{positioning}
%%%%%%%%%%%%%%%%%%%
%%%%%%% Definitions
\theoremstyle{definition}
\newtheorem{defn}{Definition}[section]
\newtheorem{fact}[defn]{Fact}
\newtheorem{prop}[defn]{Proposition}
\newtheorem{thm}[defn]{Theorem}
\newtheorem{example}[defn]{Example}

\newtheorem{corr}[defn]{Corollary}
\newtheorem{lem}[defn]{Lemma}
\newtheorem{question}[defn]{Question}
\newtheorem{remark}[defn]{Remark}
\newtheorem{observation}[defn]{Observation}
\newtheorem{claim}[defn]{Claim}
\newtheorem{subclaim}[defn]{Subclaim}
%%%%%%% Titles
\setlength{\textheight}{1.1\textheight}
%%%%%%%%%%%%%%%%%%%%%
\title[The List-distinguishing chromatic number]{The List-distinguishing chromatic number of graphs containing only small complete bigraphs}

%{Proper Distinguishing Coloring of graphs containing only small complete bigraphs}

\author{Amitayu Banerjee}
\address{E\"otv\"os Lor\'and University, Budapest, Hungary}
\email{banerjee.amitayu@gmail.com}

\date{}

\makeatletter
\@namedef{subjclassname@2020}{\textup{2020} Mathematics Subject Classification}
\makeatother
\subjclass[2020]{05C15, 05C25.}
\keywords{List-distinguishing chromatic number, Cayley graphs, Menger's theorem, Distinguishing chromatic number}
\begin{document}
\begin{abstract}
Let $G$ be a simple finite graph with maximum degree $\Delta(G)$.
The distinguishing chromatic number $\chi_{D}(G)$ of $G$ is the minimum number of colors in a proper vertex coloring of $G$ that is preserved only by the identity automorphism. Collins and Trenk (2006) proved that $\chi_{D}(G)\leq 2 \Delta(G)$ for any connected graph $G$, and that the equality holds if $G$ is the complete bigraph $K_{\Delta(G),\Delta(G)}$ or $C_{6}$. 
Inspired by Catlin's combinatorial techniques from 1978, we establish improved bounds for several classes of connected graphs that have only small complete bigraphs as induced subgraphs.
Fix $4\leq r\leq \frac{\Delta(G)}{2}$.
We prove that if $G$ is an $(n-2r+1)$-connected graph of order $n>2r$, and 
$G$ does not contain complete bigraphs $K_{r,r+1}$ and $K_{r+1,r+1}$ as induced subgraphs, then $\chi_{D}(G)\leq \chi_{D_{L}}(G)\leq 2\Delta(G)- (3\lfloor\frac{(\Delta(G)+2)}{r+1}\rfloor-4)$ where $\chi_{D_{L}}(G)$ is the list-distinguishing chromatic number of $G$ and $\chi_{D}(H)<\chi_{D_{L}}(H)$ holds in general for a connected graph $H$. 
We apply Menger's theorem to demonstrate applications of the result for graphs related to Paley graphs, Cayley graphs on Dihedral groups, and circulant Cayley graphs.
\end{abstract}

\maketitle
\section{Introduction} 
Let $K_{n}$ be a complete graph of $n$ vertices, $K_{t_{1},...,t_{m}}$ be a complete $m$-partite graph with partitions of size $t_{1},...,t_{m}$, and $C_{n}$ be a cycle of $n$ vertices.
Collins and Trenk \cite{CT2006} introduced the distinguishing chromatic number of a graph and obtained a general upper bound for it.

\begin{thm}\label{Theorem 1.1}{(Collins and Trenk \cite[Theorem 4.5]{CT2006})}
{\em Let $G$ be a connected graph. Then $\chi_{D}(G)\leq 2\Delta(G)-1$ if $G\not\in\{K_{\Delta(G),\Delta(G)},C_{6}\}$ and $\chi_{D}(G) = 2\Delta(G)$} if $G\in\{K_{\Delta(G),\Delta(G)},C_{6}\}$.    
\end{thm}

Distinguishing chromatic number of several classes of graphs has been intensively
studied in graph theory.
Laflamme and Seyffarth \cite{LS2009} proved that if $G$ is a connected bipartite graph such that $\Delta(G)\geq 3$, then 
$\chi_{D}(G)\leq 2\Delta(G)-2$ whenever $G\not\cong K_{\Delta(G)-1,\Delta(G)}, K_{\Delta(G),\Delta(G)}$.
Fijav\v{z}, Negami, and Sano \cite{FNS2011} proved that $\chi_{D}(G) \leq  5$ for every 3-connected planar graph $G \not\in \{K_{2,2,2}, C_{6} + \overline{K_{2}}\}$.
For any fixed $k \in \mathbb{N}$, Balachandran, Padinhatteer, and Spiga \cite{BPS2019} found an infinite family of Cayley graphs $G_{n}=(V(G_{n}), E(G_{n}))$ with $\chi_{D}(G_{n}) > \chi(G_{n})>k$ such that $\vert Aut(G_{n})\vert < 2k\vert V(G_{n})\vert$ where $\chi(G)$ is the chromatic number of $G$. 
In Banerjee, Gopaulsingh, and Moln\'{a}r \cite[Theorem 2.8]{BGMa2025}, we proved that if $G$ is a connected graph of order at least $3$ and $M(G)$ is the middle graph of $G$, then $\chi_{D}(M(G))$ is $\Delta(G)+1$ if $G\not\in\{C_{4}, K_{4}, C_{6}, K_{3,3}\}$, and $\Delta(G)+2$ otherwise. 
Alikhani and Soltani \cite{AS2016} proved that if $G$ is a bipartite graph of girth at least six, then   
$\chi_{D}(G)\leq \Delta(G)+1$. 
Later, Cranston \cite{Cra2018} confirmed a conjecture of Collins and Trenk by showing that a connected graph $G \not\cong C_{6}$ of girth at least five satisfies $\chi_{D}(G)\leq \Delta(G)+1$.

\begin{defn}\label{Definition 1.2}
Given a family of graphs $H_{1},..., H_{k}$, a graph $G$ is {\em $(H_{1},...,H_{k})$-free} if none of the graphs $H_{1}$,...,$H_{k}$ is an induced subgraph of $G$.    
\end{defn}

Cranston \cite{Cra2018} proved that if $G$ is a connected $(C_{3}, C_{4})$-free graph, then $\chi_{D}(G) \leq \Delta(G) + 1$ unless $G \cong C_{6}$. Recently, Brause, Kalinowski, Pil\'{s}niak, and Schiermeyer \cite{BKPS2025} showed that the upper bound of $\chi_{D}(G)$ mentioned in Theorem \ref{Theorem 1.1} can be substantially reduced if one forbids certain small graphs as induced subgraphs of $G$. 
In 2013, Ferrara et al. \cite{FGHSW2013} extended the notion of distinguishing proper coloring to a list distinguishing proper coloring. 

\begin{defn}\label{Definition 1.3}
Given an assignment $L = {L(v)}_{v\in V(G)}$ of lists of available colors to the vertices of $G$, we say that $G$ is {\em properly L-distinguishable} if there is a proper coloring $f$ of $G$ such that
$f$ is preserved only by the identity automorphism and
$f(v) \in L(v)$ for all $v\in V(G)$. 
The {\em list-distinguishing chromatic number} of $G$, denoted by $\chi_{D_{L}}(G)$, is the minimum integer $k$ such that $G$ is properly $L$-distinguishable for any list
assignment $L$ with $\vert L(v)\vert = k$ for all $v\in V(G)$.
\end{defn}

It is clear that $\chi(G)\leq \chi_{D}(G) \leq \chi_{D_{L}}(G)$. In Banerjee, Moln\'{a}r, and Gopaulsingh \cite[Section 1.1]{BMG2025}, we constructed a connected graph $G'$ such that $\chi_{D_L}(G')\neq \chi_D(G')$ and proved the following.

\begin{thm}\label{Theorem 1.4}
{(Banerjee, Moln\'{a}r, and Gopaulsingh \cite[Theorem 3.1]{BMG2025})} 
{\em Let $G$ be a connected finite graph with maximum degree $\Delta(G)$. Then $\chi_{D_{L}}(G)\leq 2\Delta(G)-1$, unless $G$ 
is $K_{\Delta(G),\Delta(G)}$ or $C_6$. In these cases, $\chi_{D_{L}}(G) = 2\Delta(G)$. } 
\end{thm}
 
Borodin and Kostochka \cite{BK1977}, Catlin \cite{Cat1978}, and Lawrence \cite{Law1978} independently improved the Brooks' upper bound for $\chi (G)$ for graphs omitting certain cliques.  
They showed that if $K_{r+1} \not\subseteq G$, where $3 \leq r \leq \Delta(G)$, then $\chi(G) \leq \frac{r}{r+1}(\Delta(G) + 2)$ by applying the following result of Lov\'{a}sz.

\begin{lem}\label{Lemma 1.5}{\em (Lov\'{a}sz; \cite{Lov1966})}
{\em If $\sum_{i=1}^{t} x_{i}\geq \Delta(G)+1-t$, then there is a partition of $V(G)$ into sets $V_{1},...,V_{t}$ such that $\Delta(G[V_{i}])\leq x_{i}$ for $1\leq i\leq t$, if $G[V_{i}]$ is the subgraph of $G$ induced by $V_{i}$.}    
\end{lem}

In \cite{BMG2025}, we applied Lemma \ref{Lemma 1.5} and followed the methods of Catlin \cite{Cat1978} to reduce the upper bound of $\chi_{D_{L}}(G)$ from Theorem \ref{Theorem 1.4} for certain graphs and proved the following result.

\begin{thm}\label{Theorem 1.6}
{(Banerjee, Moln\'{a}r, and Gopaulsingh \cite[Theorem 4.3]{BMG2025})} 
{\em Fix $6\leq t\leq \Delta(G)$. Let $G$ be a $(n-t)$-connected graph of order $n$ such that $G$ is $K_{t,t}$-free.
Then, 
\begin{center}
$\chi_{D_{L}}(G)\leq 2\Delta(G)-  (3\lfloor\frac{(\Delta(G)+1)}{t+1}\rfloor-2)$.    
\end{center}
}
\end{thm}

\subsection{Main Result} In \cite{Cat1978a}, Catlin used an intriguing combinatorial trick instead of applying Lemma \ref{Lemma 1.5} to prove that if $G$ contains no $K_{r+2}\backslash \{e\}$ as a subgraph,\footnote{We follow the notations of Catlin \cite{Cat1978a}, where $K_{r+2}\backslash \{e\}$ denotes the clique on $r+2$ vertices, minus an edge.} where $3\leq r$, then $\chi(G)\leq \frac{r}{r+1}(\Delta(G)+3)$.
In this manuscript, inspired by the methods of Catlin \cite{Cat1978a}, we reduce the upper bound of $\chi_{D_{L}}(G)$ from Theorem \ref{Theorem 1.4} for a class of graphs and prove the following result.

\begin{thm}\label{Theorem 1.7}{\em Fix $4\leq r\leq \frac{\Delta(G)}{2}$. Let $G$ be a $(n-2r+1)$-connected graph of order $n>2r$ such that
$G$ is a ($K_{r,r+1}, K_{r+1,r+1}$)-free graph.
Then, 
\begin{center}
$\chi_{D_{L}}(G)\leq 2\Delta(G)-  (3\lfloor\frac{(\Delta(G)+2)}{r+1}\rfloor-4)$.
\end{center}
}
\end{thm}

\subsection{Applications} 
Fix integers $r>6$, $1\leq k\leq 5$, and $n$ with $3r+3\le n\le r^2$. Let $a=n-2r+1$ and $b=2r-1$.
\begin{enumerate}
    \item Let $P(a)$ be the Paley graph where $a$ is a power of a prime with $a\equiv 1\pmod4$ satisfying $2r+5 \leq a \leq \min\{4r-3,\ (r-1)^2\}$.

    \item Let $D_{2m}=\langle \rho,\sigma : \rho^m=\sigma^2=1,
    \sigma\rho\sigma=\rho^{-1}\rangle$ 
be the dihedral group of order $a=2m\geq r+8$, and 
$S = (R\setminus\{e\}) \cup F$
be the generating set where $R\setminus\{e\}=\{\rho^j:1\le j\le m-1\}$ is the set of all non-trivial rotations, and 
$F\subseteq \{\rho^j\sigma:0\le j\le m-1\}$ is a collection of $m-k$ reflections.
Let $\mathrm{Cay}(D_{2m},S)$ be the Cayley graph based on $D_{2m}$ and $S$. 
     
     \item Let $C^{t}_{a}$ be the $t^{th}$-power of $C_{a}$ where $t=\max\Big\{\Big\lceil\frac{a}{r-1}\Big\rceil-1,\ \Big\lceil\frac{a-b}{2}\Big\rceil,\ \Big\lceil\frac{r+2}{2}\Big\rceil\Big\}$ and $2t\leq b$.
\end{enumerate}

Let $G_{1}\lor G_{2}$ be the join of two graphs $G_{1}$ and $G_{2}$ and let $\alpha(G)$ be the size of the maximum independent set of $G$. Define $\mathcal{G}=\bigcup_{1\leq i\leq 4}\mathcal{G}_{i}$ where

\begin{itemize}
   \item $\mathcal{G}_{1}=\Bigl\{\,K_{n-2r+1} \vee H: H \text{ is any graph on } 2r-1 \text{ vertices with }\alpha(H)\leq r-1\Bigl\}\,$,

    \item $\mathcal{G}_{2}=\Bigl\{\,P(a) \vee H: H \text{ is any graph on } 2r-1 \text{ vertices with }\alpha(H)\leq r-1\Bigl\}\,$,

    \item $\mathcal{G}_{3}=\Bigl\{\,\text{Cay}(D_{2m},S) \vee H: H \text{ is any graph on } 2r-1 \text{ vertices with }\alpha(H)\leq r-1\Bigl\}\,$, and 

    \item $\mathcal{G}_{4}=\Bigl\{\,C^{t}_{a} \vee H: H \text{ is any graph on } 2r-1 \text{ vertices with }\alpha(H)\leq r-1\Bigl\}\,$.
\end{itemize}

In Theorem \ref{Theorem 4.1} and Propositions \ref{Proposition 4.3},  \ref{Proposition 4.4}, and \ref{Proposition 4.5}, we show that for every $r>6$ and every integer $n$ with $3r+3\le n\le r^2$, any $G\in\mathcal{G}$
satisfies the following properties:

\begin{enumerate}[$(i)$]
    \item $G$ is $(n-2r+1)$-connected,
    \item $G$ has no induced subgraph isomorphic to $K_{r,r+1}$ or $K_{r+1,r+1}$, and
    \item $\Delta(G)\geq 3r+1$.
\end{enumerate}

Applying Theorem \ref{Theorem 1.7}, we can see that $\chi_{D_{L}}(G)\leq 2\Delta(G)-(9-4)=2\Delta(G)-5$ is a better bound for $\chi_{D_{L}}(G)$ than $\chi_{D_{L}}(G)\leq 2\Delta(G)-1$.

\section{Definitions and useful facts}

\begin{defn}\label{Definition 2.1}
Let $G$ be a graph with the vertex set $V(G)$ and the edge set $E(G)$.
\begin{enumerate}
\item $G$ is {\em $k$-connected} if it has at least $k+1$ vertices and removal of any $k-1$ or fewer vertices leaves a connected graph. 
The connectivity of $G$ is the largest $k$ for which $G$ is $k$-connected.
Let $\kappa(G)$ be the connectivity of the graph $G$.

\item The {\em degree} of a vertex $v \in V(G)$ in $G$, denoted by $\deg_{G}(v)$, is the number of edges that emerge from $v$.
We say $G$ is a {\em $k$-regular graph} if all vertices of $G$ have degree $k$. The degree of a regular graph $G$, denoted by $d(G)$, is the same as the maximum degree $\Delta(G)$ and the minimum degree $\delta(G)$.

\item $G$ is an {\em edge-transitive graph} if given any two edges $e_{1}$ and $e_{2}$ of $G$, there is an automorphism of $G$ that maps $e_{1}$ to $e_{2}$. 
%A graph that is both edge-transitive and vertex-transitive is called a {\em symmetric} graph.

\item An {\em independent set} of $G$ is a set of vertices of $G$, no two of which are adjacent vertices. 
%Let $\alpha(G)$ be the size of the maximum independent set of $G$.
\end{enumerate}
\end{defn}

Let $x,y \in V(G)$ be two distinct non-adjacent vertices. Menger's theorem states that the minimum number of vertices, distinct from $x$ and $y$,  whose removal disconnects $x$ and $y$ 
is equal to the maximum number of pairwise internally vertex-disjoint $x$--$y$ paths in $G$.
Consequently, $G$ is $k$-connected if every pair of vertices has at least $k$ internally vertex disjoint paths in between.

\begin{lem}{(Hoffman; \cite[Theorem 1]{Hae2021})}\label{Lemma 2.2}
{\em The independence number $\alpha(G)$ for a $k$-regular graph $G$ on $n$ vertices with the least eigenvalue $\lambda_{min}$ is at most
$n\cdot\frac{-\lambda_{\min}}{k-\lambda_{\min}}$.}   
\end{lem}

\begin{defn}[Undirected Cayley graph]\label{Definition 2.3}
Let $\Gamma$ be a finite group and let $S \subseteq \Gamma \setminus \{e\}$ 
be an inverse-closed subset of $\Gamma \setminus \{e\}$ i.e., $S = S^{-1}$ , where 
$S^{-1}:= \{s^{-1}: s \in S\}$.
The \emph{undirected Cayley graph} Cay$(\Gamma, S)$ is the graph with a set of vertices $\Gamma$, and the vertices $u$ and $v$ are adjacent in Cay$(\Gamma, S)$ if and only if $uv^{-1} \in S$. 
\end{defn}

It is known that Cay$(\Gamma, S)$ is a $|S|$-regular graph, and Cay$(\Gamma, S)$ is
connected if and only if $S$ is a generating set of $G$.

\begin{defn}{($k$-th power of a cycle)}\label{Definition 2.4}
Let $C_{n}$ be the cycle on $n$ vertices.  
For a positive integer $k$, the \emph{$k$-th power of the cycle} $C_{n}^{k}$ is the graph with vertex set $\mathbb{Z}_{n}$ in which two distinct vertices $x,y\in \mathbb{Z}_{n}$ are adjacent whenever $\min\{\,|x-y|,\; n-|x-y|\,\}\le k$.
Equivalently, $C_{n}^{k}$ is the undirected Cayley graph
$\mathrm{Cay}\big(\mathbb{Z}_{n},\;\{\pm 1,\pm 2,\dots,\pm k\}\big)$.
\end{defn}

\begin{lem}{(Ichishima--Muntaner-Batle--Takahashi;\cite[Lemma 5]{IAT2025})}\label{Lemma 2.5}
{\em For any two integers $n$ and $k$ with $n \geq k + 2$,
$\alpha(C^{k}_{n})=\lfloor \frac{n}{k+1} \rfloor$ where $k\in [2,\lfloor \frac{n}{2}\rfloor]$.
%, and $\alpha(C^{k}_{n})$ is the size of the maximum independent set of $C^{k}_{n}$.
}    
\end{lem}

\begin{defn}{(Paley graph)}\label{Definition 2.6}
Let $q$ be a prime power with $q \equiv 1 \pmod{4}$. The \emph{Paley graph} $P(q)$ is the undirected Cayley graph of the additive group $(\mathbb{F}_{q},+)$ with generating set equal to the set of non-zero quadratic residues in $\mathbb{F}_{q}$. The vertex set of $P(q)$ is $\mathbb{F}_{q}$ and two distinct vertices $x,y \in \mathbb{F}_{q}$ are adjacent if and only if $x-y$ is a non-zero quadratic residue in $\mathbb{F}_{q}^{\times}$.
\end{defn}

\begin{lem}{(Watkins; \cite[Corollary 1A]{Wat1970})}\label{Lemma 2.7}
{\em If $G$ is a connected edge-transitive graph, then $\kappa(G)=\delta(G)$}. 
\end{lem}

The Paley graph $P(q)$ is a connected, edge transitive graph. By Lemma \ref{Lemma 2.7}, $\kappa(P(q))=\delta(P(q))$.

\section{The proof of Theorem 1.7}
Suppose $G$ is a $(n-2r+1)$-connected graph of order $n$ such that
$G$ does not contain $K_{r,r+1}$ and $K_{r+1,r+1}$ as induced subgraphs.

\subsection{Defining $t$ and $h_{i}$ for each $1\leq i\leq t$} We define 
\begin{itemize}
    \item $h=\Delta(G)$,
    \item $t=\lfloor \frac{h+2}{r+1}\rfloor$, 
    \item $h_{i}=r$ for all $1\leq i\leq t-1$, and 
    \item $h_{t}$ as an integer that satisfies $r\leq h_{t}=h+2-t(r+1)+r\leq 2r$.
\end{itemize}
Thus, $h_{1}+...+h_{t-1}+h_{t}=(t-1)r+h+2-t(r+1)+r=h-t+2$.

\subsection{Defining a partition $(X_{1},...,X_{t})$ of $V(G)$} For a subset $X\subseteq V(G)$ of vertices of a graph $G$, let $G[X]$ denote the subgraph of $G$ induced by $X$. If $(X_{1},...,X_{t})$ is a partition of $V(G)$, then we define the integer-valued function $f$ as follows:
\begin{center}
$f(X_{1},...,X_{t})=h_{1}\vert X_{1}\vert + h_{2}\vert X_{2}\vert +... + h_{t}\vert X_{t}\vert-(\vert E(G[X_{1}])\vert +...+ \vert E(G[X_{t}])\vert)$.
\end{center}

Consider the partition $(X_{1},...,X_{t})$ of $V(G)$ that
\begin{enumerate}
    \item maximizes $f(X_{1},...,X_{t})$;
    \item minimizes the total number of complete bipartite induced subgraphs $K_{r,r}$ in the $G[X_{i}]$'s in case $h_{i}=r$ and $(1)$ holds.
\end{enumerate}

Using the methods used in \cite{Cat1978a}, we obtain $\Delta(G[X_{i}])\leq h_{i}$ for all $1\leq i\leq t$. 
For the reader's convenience, we write the details from \cite{Cat1978a}.
For each $2\leq i\leq t$, we have
\[
0 \leq f(X_1, \ldots, X_t) - f(X_1 - x, X_2, \ldots, X_i + x, \ldots, X_t),
\]
which expands as
\[
0\leq h_1|X_1| - h_1(|X_1| - 1) + h_i|X_i| - h_i(|X_i| + 1) - |E(X_1)| + |E(X_1 - x)| - |E(X_i)| + |E(X_i + x)|.
\]
Simplifying, we obtain 
$0\leq h_1 - h_i - \deg_{G[X_1]}(x) + \deg_{G[X_i + x]}(x).$
Hence, it follows that
\begin{center}
$\deg_{G[X_1]}(x) \leq h_1 - h_i + \deg_{G[X_i + x]}(x)$ for each $2\leq i\leq t$.     
\end{center}

Clearly, $\deg_{G[X_{1}]}(x)\leq \deg_{G[X_{1}+x]}(x)=h_{1}-h_{1}+\deg_{G[X_{1}+x]}(x)$.
Considering the summation of both sides of this system of inequalities by letting $i$ run from $2$ to $t$, we obtain
\[
t \cdot \deg_{G[X_1]}(x) \leq t h_1 - \sum_{i=1}^{t} h_i + \sum_{i=1}^{t} \deg_{G[X_i + x]}(x).
\]
We note that $h=\Delta(G)$,
$\sum_{i=1}^{t} \deg_{G[X_i + x]}(x) = \deg_G(x)$,
and \( \sum_{i=1}^{t} h_i = h - t + 2 \). Thus, 
\[
t \cdot \deg_{G[X_1]}(x) \leq t h_1 - (h - t + 2) + \deg_G(x) \leq t h_1 + t - 2 \implies 
\deg_{G[X_1]}(x) \leq h_1 + \frac{t - 2}{t}.
\]
Since both \( \deg_{G[X_1]}(x) \) and \( h_1 \) are integers, it follows that
$
\deg_{G[X_1]}(x) \leq h_1.
$
A similar argument applies for any vertex \( x \in X_i \) with \( i \leq t \), yields
$
\deg_{G[X_i]}(x) \leq h_i.
$

\subsection{Main proof}

\begin{claim}\label{claim 3.1}
{\em If $h_{t}>r$, then $\chi_{D_{L}}(G[X_{t}])\leq 2h_{t}-1$.} 
\end{claim}

\begin{proof}

\noindent \textsc{Case 1:}  Suppose $\vert V_{G[X_{t}]}\vert< 2r$.

As coloring the vertices with distinct colors from their respective lists of size $\vert V_{G[X_{t}]}\vert$ yields a proper distinguishing coloring we have
\[\chi_{D_{L}}(G[X_{t}])\leq \vert V_{G[X_{t}]}\vert \leq 2r-1\leq 2(h_{t}-1)-1< 2h_{t}-1.\] 

\noindent \textsc{Case 2:} Suppose $\vert V_{G[X_{t}]}\vert\geq 2r$. Since $G$ is $(n-2r+1)$-connected, we have that $G[X_{t}]$ is connected.
We note that $\Delta(G[X_{t}])\leq h_{t}$.

\vspace{2mm}
\begin{enumerate}
    \item[] \textsc{Subcase 2.1:} Let $\Delta(G[X_{t}])< h_{t}$.
Then, by Theorem \ref{Theorem 1.4}, we have
\[\chi_{D_{L}}(G[X_{t}])\leq 2(\Delta (G[X_{t}]))\leq 2(h_{t}-1)=2h_{t}-2.\]
    %\vspace{2mm}
    \item[] \textsc{Subcase 2.2:}
    Let $\Delta(G[X_{t}])= h_{t}>r$.
Since $G$ contains no $K_{r+1,r+1}$ as an induced subgraph (and hence no $K_{h_{t}, h_{t}}$), neither does $G[X_{t}]$. 
Thus, $G[X_{t}]\neq K_{\Delta(G[X_{t}]), \Delta(G[X_{t}])}$. 
Moreover, $G[X_{t}]\neq C_{6}$ as $\vert V_{G[X_{t}]}\vert\geq 2r\geq 8$.
Thus, by Theorem \ref{Theorem 1.4}, we have
\[\chi_{D_{L}}(G[X_{t}])\leq 2\Delta(G(X_{t}))-1=2h_{t}-1.\]
\end{enumerate} 
This finishes the proof of Claim \ref{claim 3.1}.
\end{proof}

\begin{claim}\label{claim 3.2}
{\em $\chi_{D_{L}}(G[X_{i}])\leq 2h_{i}-1$ for all $1\leq i\leq t$.}
\end{claim}

\begin{proof}
In view of Claim \ref{claim 3.1}, we may assume that $h_{i}=r$ for all $1\leq i\leq t$.
We fix $1\leq i\leq t$.
\noindent \textsc{Case 1:}  If $\vert V_{G[X_{i}]}\vert< 2r$, then
$\chi_{D_{L}}(G[X_{i}])\leq \vert V_{G[X_{i}]}\vert \leq 2r-1= 2h_{i}-1.$ 

\noindent \textsc{Case 2:} If $\vert V_{G[X_{i}]}\vert\geq 2r$, then since $G$ is $(n-2r+1)$-connected, we have that $G[X_{i}]$ is connected.
We recall that $\Delta(G[X_{i}])\leq h_{i}$.

\vspace{2mm}
\begin{enumerate}
    \item[] \textsc{Subcase 2.1:} Let $\Delta(G[X_{i}])< h_{i}$.
Then, by Theorem \ref{Theorem 1.4}, we have
\[\chi_{D_{L}}(G[X_{i}])\leq 2(\Delta (G[X_{i}]))\leq 2(h_{i}-1)=2h_{i}-2.\]
    %\vspace{2mm}
    \item[] \textsc{Subcase 2.2:}
    Let $\Delta(G[X_{i}])= h_{i}=r$.

\begin{subclaim}\label{Subclaim 3.3}
{\em There is no induced subgraph $K_{\Delta(G[X_{i}]), \Delta(G[X_{i}])}$ in $G[X_{i}]$.}    
\end{subclaim}

\begin{proof}
For the sake of contradiction, suppose that $G[X_{i}]$ contains an induced subgraph $C_{0}=K_{\Delta(G[X_{i}]), \Delta(G[X_{i}])}=K_{h_{i},h_{i}}=K_{r,r}$. Let $x_{0}\in C_{0}$. 
\vspace{2mm}

\textbf{Step 1:} 
We note that $x_{0}$ is adjacent to  $\Delta(G[X_{i}])=h_{i}$ vertices in $C_{0}=K_{h_{i}, h_{i}}$.
{\em We claim there exists $1\leq j\leq t$ such that $j\neq i$, and $x_{0}$ is adjacent to $h_{j}$ or fewer vertices in $X_{j}$.} 
Otherwise, if $x_{0}$ is adjacent to at least $h_{j}+1$ vertices in each $X_{j}$ for $j\neq i$, then 

\begin{center}
$\deg_{G}(x_{0})\geq h_{i}+\sum_{j\neq i, 1\leq j\leq t}(h_{j}+1)=\sum_{1\leq p\leq t}h_{p}+(t-1)=(h-t+2)+t-1=h+1>h$,     
\end{center}

since $\sum_{1\leq p\leq t}h_{p}=h-t+2$, which is a contradiction to the assumption that $h=\Delta(G)$. We fix such a $j$. 
\vspace{2mm}

\textbf{Step 2:} Suppose, we move $x_{0}$ from $X_{i}$ to $X_{j}$. Then the maximality of $f$ is preserved as
$f(X_{1},..., X_{t}) = f(X_{1},.., X_{i}-x_{0},..., X_{j}+x_{0},.., X_{t})$.
\vspace{2mm}

\textbf{Step 3:} We obtain the desired contradiction to the assumption that $G[X_{i}]$ contains $C_{0}$. To avoid violating condition (2), i.e. the partition $(X_{1},..., X_{t})$ minimizes the total number of induced subgraph $K_{r,r}$ in the $G[X_{i}]$'s in case $h_{i}=r$ and $f(X_{1},..., X_{t})$ is maximum, 
the destruction of $C_{0}=K_{r,r}$ in $G[X_{i}\backslash \{x_{0}\}]$ is followed by the formation of another induced subgraph $K_{r,r}$ in $G[X_{j}\cup \{x_{0}\}]$, say $C_{1}$. 
Then the complete bipartite graph $C_{0}\backslash\{x_{0}\}=K_{r-1,r}$ is left behind in $G[X_{i}\backslash \{x_{0}\}]$ (See Fig. \ref{Figure 1}).

\newcommand{\rsize}{5} % number of vertices in each bipartite part

\begin{figure}[h]
\centering

\begin{tikzpicture}[scale=0.85, every node/.style={inner sep=1.5pt}]

% outer rectangle
\def\W{12}   % total width
\def\H{4.5}  % total height
\def\sqW{\W/4} % width of each small square

\draw[thick] (0,0) rectangle (\W,\H);
\foreach \i in {1,2,3} {
  \draw[thick] (\i*\sqW,0) -- (\i*\sqW,\H);
}

%%%%%%%%%%%%%%%%%%%%%%%%%%%%%%%%%%%%%%%
% --- inner rounded square (C0) inside first square ---
\def\inset{0.3}
\coordinate (A0) at (\inset,\inset);
\coordinate (B0) at (\sqW-\inset,\H-\inset);

\draw[thick,rounded corners=12pt] (A0) rectangle (B0);

% label C0 ABOVE the rounded square (outside)
\node[font=\small, anchor=south, yshift=-6mm, xshift=+7mm] 
  at ($ (A0)!0.5!(B0) $ |- B0) {$C_{0}$};

% coordinates for bipartite graph inside C0
\pgfmathsetmacro{\leftx}{\inset+0.25*(\sqW-2*\inset)}   
\pgfmathsetmacro{\rightx}{\inset+0.75*(\sqW-2*\inset)}
\pgfmathsetmacro{\topY}{\H-\inset-0.2}   % full height again
\pgfmathsetmacro{\bottomY}{\inset+0.2}

\pgfmathtruncatemacro{\denom}{\rsize-1}
\ifnum\denom<1
  \pgfmathsetmacro{\denom}{1}
\fi

\foreach \i in {1,...,\rsize} {
  \pgfmathsetmacro{\y}{\bottomY + (\topY-\bottomY)*(\rsize-\i)/\denom}
  \node[draw,circle,fill=white,inner sep=1pt] (L\i) at (\leftx,\y) {};
  \node[draw,circle,fill=white,inner sep=1pt] (R\i) at (\rightx,\y) {};
}
\node[below left=1.5pt and -2pt] at (L1.south west) {$x_{0}$};

\foreach \i in {1,...,\rsize} {
  \foreach \j in {1,...,\rsize} {
    \draw (L\i) -- (R\j);
  }
}

%%%%%%%%%%%%%%%%%%%%%%%%%%%%%%%%%%%
% --- inner rounded square (C1) inside third square ---
\coordinate (A1) at (2*\sqW+\inset,\inset);
\coordinate (B1) at (3*\sqW-\inset,\H-\inset);

\draw[thick,rounded corners=12pt] (A1) rectangle (B1);

% label C1 ABOVE the rounded square (outside)
\node[font=\small, anchor=south, yshift=-6mm, xshift=7mm] 
  at ($ (A1)!0.5!(B1) $ |- B1) {$C_{1}$};

% coordinates for bipartite graph inside C1
\pgfmathsetmacro{\leftxC}{2*\sqW+\inset+0.25*(\sqW-2*\inset)}   
\pgfmathsetmacro{\rightxC}{2*\sqW+\inset+0.75*(\sqW-2*\inset)}
\pgfmathsetmacro{\topYC}{\H-\inset-0.2}   % full height again
\pgfmathsetmacro{\bottomYC}{\inset+0.2}

\foreach \i in {1,...,\rsize} {
  \pgfmathsetmacro{\y}{\bottomYC + (\topYC-\bottomYC)*(\rsize-\i)/\denom}
  \node[draw,circle,fill=white,inner sep=1pt] (Lc\i) at (\leftxC,\y) {};
  \node[draw,circle,fill=white,inner sep=1pt] (Rc\i) at (\rightxC,\y) {};
}
\node[below left=2pt and -2pt] at (Lc1.south west) {$x_{0}$};

\foreach \i in {1,...,\rsize} {
  \foreach \j in {1,...,\rsize} {
    \draw (Lc\i) -- (Rc\j);
  }
}
%%%%%%%%%%%%%%%%%%%%%%%%%%%%%
% label X_{i} and X_{j}
\node[font=\small, anchor=south, yshift=-25mm, xshift=0mm] 
  at ($ (A1)!0.5!(B1) $ |- B1) {$G[X_{j}\cup \{x_{0}\}]$};
  
\node[font=\small, anchor=south, yshift=-25mm, xshift=0mm] 
  at ($ (A0)!0.5!(B0) $ |- B0) {$G[X_{i}\backslash \{x_{0}\}]$};
%%%%%%%%%%%%%%%%%%%%%%%%%%%%%%%%%%%%%
% curved arrow from x_{0} in C0 to x_{0} in C1
\draw[->,thick,bend left=12] (L1) to (Lc1);

\end{tikzpicture}
\caption{\em After vertex $x_{0}$ is moved from $X_{i}$ to $X_{j}$, $C_{0}\backslash \{x_{0}\}=K_{r-1,r}$ is left in $G[X_{i}\backslash\{x_{0}\}]$, and $C_{1}=K_{r,r}$ is formed in $G[X_{j}\cup \{x_{0}\}]$ when $r=5$.}
\label{Figure 1}
\end{figure}

We repeat this process by picking a vertex $x_{1}\neq x_{0}$ in $C_{1}$ and removing it from $G[X_{j}\cup \{x_{0}\}]$ to create another induced subgraph $K_{r,r}$ in $G[X_{k}\cup \{x_{1}\}]$ for some $1\leq k\leq t$, $k\neq j$, say $C_{2}$. This leaves behind a complete bipartite graph $K_{r-1,r}$ in $G[X_{j}\backslash \{x_{1}\}]$. 
Continuing in this fashion, we will obtain a sequence $\{C_{i}\}_{i\geq 0}$.
Since $G$ is finite, there will be a vertex $x_{m-1}$ that will move to $X_{l}$ for some $1\leq l\leq t$, where it will form an induced subgraph $C_{m}=K_{r,r}$ (a member of the sequence $\{C_{i}\}_{i\geq 0}$) in $G[X_{l}\cup \{x_{m-1}\}]$ and $C_{m}$ overlaps with another member $C_{k}$ of the sequence $\{C_{i}\}_{i\geq 0}$ for some $k<m-1$. 
Without loss of generality, we choose $x_{m-1}\neq x_{k}$ from $C_{m-1}$.

\vspace{2mm}

Then $C_{m}-x_{m-1} \cong C_{k}-x_{k}$, a complete bipartite graph $K_{r-1,r}$ in $G[X_{l}]$, is left behind as some vertex $x_{k}$
in the sequence
$x_{0}, x_{1},..., x_{m-2}$
was moved out.
Finally, $x_{m-1}, x_{k}$, and $C_{m}-x_{m-1}\cong C_{k}-x_{k}=K_{r-1,r}$ forms an induced complete bipartite subgraph $K_{r+1,r}$ of $G$ which contradicts the hypothesis of the Theorem (see Fig. \ref{Figure 2}).
\vspace{2mm}

Thus, there is no induced subgraph $K_{\Delta(G[X_{i}]), \Delta(G[X_{i}])}$ in $G[X_{i}]$.
\end{proof}


\newcommand{\rsize}{5} % this will make A size = r-1 = 4, B size = 5



\begin{figure}[h]
\centering


\begin{tikzpicture}[scale=0.8]

% --- even smaller square container, moved down ---
\def\sqsize{4} % square size
\def\yshift{0.3} % shift down
\draw[thick] (0.7,0.1) rectangle (3.3,3.3);

% label outside, top-right, slightly lower
\node[anchor=north west] at (3.5, 3.5) {$C_{m}-x_{m-1} \cong C_{k}-x_{k}=K_{r-1,r}$};

% --- bipartition vertices inside square ---
% Partition A (4 vertices) on the left
\foreach \i in {1,...,4} {
  \fill (1,\sqsize+\yshift-0.5-\i*0.7) circle (2pt);
}

% Partition B (5 vertices) on the right
\foreach \j in {1,...,5} {
  \fill (\sqsize-1,\sqsize+\yshift-0.5-\j*0.7) circle (2pt);
}

% --- edges of K_{5,4} ---
\foreach \i in {1,...,4} {
  \foreach \j in {1,...,5} {
    \draw (1,\sqsize+\yshift-0.5-\i*0.7) -- (\sqsize-1,\sqsize+\yshift-0.5-\j*0.7);
  }
}

% --- vertices x_k and x_{m-1} above the 4-vertex partition, as black circles ---
\fill (1,\sqsize+\yshift-0.5) circle (2pt) node[left=1pt] {$x_{k}$};
\fill (1,\sqsize+\yshift+0.2) circle (2pt) node[left=1pt] {$x_{m-1}$};

% connect x_k and x_{m-1} to all vertices of partition B
\foreach \j in {1,...,5} {
  \draw (1,\sqsize+\yshift-0.5) -- (\sqsize-1,\sqsize+\yshift-0.5-\j*0.7);
  \draw (1,\sqsize+\yshift+0.2) -- (\sqsize-1,\sqsize+\yshift-0.5-\j*0.7);
}
\end{tikzpicture}
\caption{\em Formation of the induced subgraph $K_{r+1,r}$ of $G$ when $r=5$.}
\label{Figure 2}
\end{figure}

By Subclaim \ref{Subclaim 3.3}, we have $G[X_{i}]\neq K_{\Delta(G[X_{i}]), \Delta(G[X_{i}])}$. 
Moreover, $G[X_{i}]\neq C_{6}$ as $\vert V_{G[X_{i}]}\vert\geq 2r\geq 8$.
By Theorem \ref{Theorem 1.4}, we have
\[\chi_{D_{L}}(G[X_{i}])\leq 2\Delta(G(X_{i}))-1=2h_{i}-1.\]
\end{enumerate}
This finishes the proof of Claim \ref{claim 3.2}.
\end{proof}


\begin{defn}\label{Definition 3.4}
We denote init$_{k}(L)$ by the first $k$ elements of a list $L$.    
\end{defn}

In \cite[Claim 4.4]{BMG2025}, we proved 
$\chi_{D_{L}}(G)\leq \sum_{i=1}^{t} \chi_{D_{L}}(G[X_{i}])$. For the reader's convenience, we provide the details here as well.

\begin{claim}{(cf. \cite[Claim 4.4]{BMG2025})}\label{claim 3.5}
$\chi_{D_{L}}(G)\leq \sum_{i=1}^{t} \chi_{D_{L}}(G[X_{i}])$.
\end{claim}

\begin{proof}
Let $k_{i}=\chi_{D_{L}}(G[X_{i}])$ for each $1\leq i\leq t$. Assign any list $L(v)$ to each vertex $v \in V(G)$ such that $\vert L(v)\vert= \sum_{i=1}^{t}k_{i}$. We define a coloring $f$ for the vertices of $G$ as follows:

\begin{enumerate}
    \item Let $L^{1}(v):=$ init$_{k_{1}}(L(v))$ for each vertex $v\in V(G)$. 
    \vspace{2mm}
    
    Let $L^{1}=\{L^{1}(v)\}_{v\in V_{G[X_{1}]}}$. Since $G[X_{1}]$ is properly $L^{1}$-distinguishable (say by coloring $f^1$), we define $f\restriction V_{G[X_{1}]}:=f^{1}$.
    \vspace{2mm}
    
    \item For any $1<i\leq t$, let $L^{i}(v):=$ init$_{k_{i}}(L(v)\backslash \sum_{k= 1}^{i-1} L^{k}(v)) $ for each vertex $v\in V(G)$. 
    \vspace{2mm}
    
    Let $L^{i}=\{L^{i}(v)\}_{v\in V_{G[X_{i}]}}$. Since $G[X_{i}]$ is properly $L^{i}$-distinguishable (say by coloring $f^i$), we define $f\restriction V_{G[V_{i}]}:=f^{i}$.
\end{enumerate}  

If $\phi\in Aut(G)$ preserves $f$, then the range of $G[X_{i}]$ with respect to $\phi$ is $G[X_{i}]$, for all $1\leq i\leq t$. Thus, $\phi\restriction G[X_{i}]$ is an automorphism of $G[X_{i}]$ preserving $f$. Since the above defined colorings of $G[X_{i}]$'s are distinguishing, we have that $\phi\restriction G[X_{i}]$ is a trivial automorphism for all $1\leq i\leq t$. Thus, $f$ is a distinguishing coloring of $G$. 
To show $f$ is a proper coloring of $G$, we pick any $x,y\in V(G)$ such that $\{x,y\}\in E_{G}$. If $x,y\in V_{G[X_{i}]}$ for some $1\leq i\leq t$, then $f(x)\neq f(y)$, as $f\restriction (V_{G[X_{i}]})$ is a proper coloring. If $x\in V_{G[X_{i}]}$ and $y\in V_{G[X_{j}]}$ such that $i\neq j$, then $f(x)\neq f(y)$, as $f$ is defined in a way such that the sets of colors used to color $G[X_{i}]$ and $G[X_{j}]$ are different.
\end{proof} 

By applying Claims 
\ref{claim 3.2} and 
\ref{claim 3.5}, we get,
\begin{equation*}
\begin{aligned}
 \chi_{D_{L}}(G) \leq \sum_{i=1}^{t} \chi_{D_{L}}(G[X_{i}])
       \leq \sum_{i=1}^{t}(2h_{i}-1)
       = 2\sum_{i=1}^{t} h_{i}-t
       =2(h-t+2)-t 
      & = 2\Delta(G)- (3t-4).
\end{aligned}
\end{equation*}
%%%%%%%%%%%%%%%%%%%%%%%%%%%%
\section{Applications}

Throughout the section, we fix integers $r>6$ and $n$ with $3r+3\le n\le r^2$. We define $a=n-2r+1$ and $b=2r-1$. 

\begin{thm}\label{Theorem 4.1}
{\em Let $X$ be any graph on $a=n-2r+1$ vertices satisfying the following properties:
\begin{enumerate}[(i)]
    \item $\alpha(X)\le r-1$, 
    \item $\kappa(X)\ge a-b$, and
    \item $\Delta(X)\ge r+2$.
\end{enumerate}

Let $H$ be any graph on $b$ vertices with $\alpha(H)\le r-1$. Then $G=X\vee H$ satisfies the following:
\begin{enumerate}
    \item $G$ is $(n-2r+1)$-connected,
    \item $G$ contains no induced subgraph isomorphic to $K_{r,r+1}$ or $K_{r+1,r+1}$,
    \item $\Delta(G)\ge 3r+1$. Thus, $\displaystyle \frac{\Delta(G) + 2}{r + 1} \ge 3$.
\end{enumerate}
Consequently, by Theorem \ref{Theorem 1.7}, we obtain $\chi_{D_{L}}(G)\leq 2\Delta(G)-5$ which is a better bound for $\chi_{D_{L}}(G)$ than $\chi_{D_{L}}(G)\leq 2\Delta(G)-1$.
}
\end{thm}

\begin{proof}
The bipartitions of an induced subgraph $K_{r,r+1}$ or $K_{r+1,r+1}$ are independent sets, so the bounds $\alpha(X)\leq r-1$ and $\alpha(H)\le r-1$ prevent $G$ from containing $K_{r,r+1}$ and $K_{r+1,r+1}$ as induced subgraphs. 
Since $G=X\lor H$, we have 
$\kappa(G)=\min\{\kappa(X)+b,\kappa(H)+a,\delta(G)\}$.
Since $\kappa(X)\ge a-b$ by assumption, we have $\kappa(X)+b\geq a$. Moreover, $\delta(G)=\delta(X\lor H)=min\{\delta(X)+b,\delta(H)+a\}\geq a$ since $\delta(X)\ge \kappa(X)$. Thus, $\kappa(G)\geq a$.
Consequently, $G$ is $a$-connected.
Finally, if $x\in V(X)$ such that $\deg_X(x)\ge r+2$ (possible by our assumption $\Delta(X)\geq r+2$), then
$\deg_G(x)=\deg_X(x)+b\geq (r+2)+(2r-1)= 3r+1$. Thus, $\Delta(G)\ge 3r+1$.
\end{proof}

\begin{example}\label{Example 4.2}
The following shows that $X=K_{a}$ satisfies $(i)-(iii)$ of Theorem \ref{Theorem 4.1}.

\begin{enumerate}[$(i)$]
    \item Since $\alpha(X) = 1$ and $r>6$, we have $\alpha(X)< r-1$.
    
    \item Since $b=2r-1> 11$, we have $\kappa(X) = a-1> a-b$. 
    
    \item Clearly, $\Delta(X) = a-1=n-2r\geq (3r+3)-2r=r+3> r+2$.
\end{enumerate}
\end{example}
%%%%%%%%%%%%%%%%%%%%%%%%%%%%%%%%
\subsection{Paley graph}

\begin{prop}\label{Proposition 4.3}
{\em Let $a$ be a power of a prime with
$a\equiv1\pmod4$ satisfying 
$2r+5 \leq a \leq \min\{4r-3,\ (r-1)^2\}$.
Then the Paley graph $X=P(a)$ satisfies $(i)-(iii)$ of Theorem \ref{Theorem 4.1}.}
\end{prop}

\begin{proof}
The graph $P(a)$ is $(\frac{a-1}{2})$-regular on $a$ vertices with the least eigenvalue $\frac{-1-\sqrt a}{2}$. Thus,
\[
\alpha(P(a))
\le a\cdot\frac{-(\frac{-1-\sqrt a}{2})}
{\frac{a-1}{2} - (\frac{-1-\sqrt a}{2})}
= a\cdot\frac{1+\sqrt a}{a+\sqrt a}
= \sqrt a
\]

by Lemma \ref{Lemma 2.2}.  As $a\le (r-1)^2$ by assumption, we obtain $\alpha(X)\le r-1$.
Now, $X$ is connected, edge-transitive, and $(\frac{a-1}{2})$-regular.  By Lemma \ref{Lemma 2.7}, $\kappa(X)=\delta(X)=\frac{a-1}{2}$. Since $4r-3\geq a$ and $b=2r-1$, we obtain $\kappa(X)\geq a-b$.
Since $X$ is $(\frac{a-1}{2})$-regular and $a\geq 2r+5$, we have $\Delta(X)=\frac{a-1}{2}\geq r+2$. 
\end{proof}

\subsection{Cayley graph of the dihedral group}

\begin{prop}\label{Proposition 4.4}
{\em Suppose $a=2m\geq r+8$ is an even integer and $1\leq k\leq 5$ is any integer. 
%So $a=2m$ for some $m\in\mathbb{N}$. 
Let $D_{2m}=\langle \rho,\sigma : \rho^m=\sigma^2=1,\ \sigma\rho\sigma=\rho^{-1}\rangle$ 
be the dihedral group of order $2m$, $R=\{\rho^j:0\le j\le m-1\}$, $R\sigma=\{\rho^j\sigma:0\le j\le m-1\}$, and $F\subseteq R\sigma$ satisfy $|F|=m-k$.
If $S \;=\; (R\setminus\{e\}) \;\cup\; F$ is a generating set,
then $X=\mathrm{Cay}(D_{2m},S)$ satisfies $(i)-(iii)$ of Theorem \ref{Theorem 4.1}.}
\end{prop}

\begin{proof}
%We only show when $k=5$. The rest of the cases are similar.
Let $d=|S|=(m-1)+(m-k)=2m-(k+1)$. Then $X$ is a $d$-regular graph.

\begin{enumerate}[(i)]
    \item We write $D_{2m}=R\cup R\sigma$ as the disjoint union of cosets $R$ and $R\sigma$.
    Now, each of $R$ and $R\sigma$ induces a clique in $X$.
Take two distinct rotations \(\rho^i,\rho^j\in R\) with \(i< j\). Then
$
(\rho^i)^{-1}\rho^j=\rho^{j-i}\in R\setminus\{e\}\subseteq S,
$
so \(\rho^i\) and \(\rho^j\) are adjacent in $X$.
Similarly, take two distinct reflections \(\rho^i\sigma,\rho^j\sigma\in R\sigma\) with \(i> j\). 
We note that $\sigma^{2}=\{e\}$, so $\sigma^{-1}=\sigma$ and $\sigma\rho^{k}\sigma=\rho^{-k}$ for all integers $k\geq 1$.
Then
\[
(\rho^i\sigma)^{-1}(\rho^j\sigma)
= (\sigma^{-1}\rho^{-i})(\rho^j\sigma) 
= \sigma\rho^{-i}\rho^j\sigma
= \sigma\rho^{\,j-i}\sigma
= \rho^{-(j-i)}=\rho^{\,i-j}\in S.
\]

So, \(\rho^i\sigma\) and \(\rho^j\sigma\) are adjacent in $X$.
Thus, any independent set in $X$ can contain at most one rotation and at most one reflection. Hence $\alpha(X)\leq 2< r-1$ since $r> 6$.
    
\item We show $\kappa(X)\geq 2m-(2k+1)$. By Menger's theorem, it suffices to show that for any two distinct vertices
$u,v\in V(X)$, there are $2m-(2k+1)$ internally vertex-disjoint $u$--$v$ paths.
By the arguments in (i), $R$ and $R\sigma$ induces a clique $K_m$ in $X$. Moreover, we can see that each rotation has exactly $m-k$ neighbours in $R\sigma$ and each reflection has exactly $m-k$ neighbours in $R$.
Pick $u,v\in V(X)$.

\medskip\noindent\textbf{Case A: $u,v\in R$ are distinct rotations.}  
Write $u=\rho^i$, $v=\rho^j$, $i\ne j$. Since $R$ induces a clique $K_{m}$, there is an edge $uv$ joining $u$ and $v$. Let $P_{0}=\{uv\}$. 
Let $R\setminus\{u,v\}=\{\rho^{r_1},\dots,\rho^{r_{m-2}}\}$. For each $\rho^{r_t}\in R\backslash\{u,v\}$ consider the path
\begin{center}
$P_t= \{u\rho^{r_t}, \rho^{r_t}v\}$.    
\end{center}

We construct $m-2k$ disjoint paths each using a distinct reflection as internal vertex.
We note that the set of reflections
adjacent to $u$ is a translate $\rho^i F$, and those adjacent to $v$ is $\rho^j F$ for some $i$ and $j$. Since $|F|=m-k$, we have $\vert N_{X}(u)\cap N_{X}(v)\cap R\sigma\vert =\vert\rho^i F\cap \rho^j F\vert\geq m-2k$.\footnote{$N_{G}(x)$ denotes the open-neighborhood of $x$ in $G$.}
Therefore, there are at least $m-2k$ reflections adjacent to both $u$ and $v$.
For each such reflection $s\in \rho^i F\cap\rho^j F$, consider the path
\[
Q_s= \{us, sv\}.
\]

Then $\{P_{0}\}\cup\{P_{t}:\rho^{r_{t}}\in R\backslash\{u,v\}\}\cup \{Q_{s}:s\in \rho^i F\cap\rho^j F, N_{X}(u)=\rho^{i}F, N_{X}(v)=\rho^{j}F\}$ is a collection of $(m-1)+(m-2k)=2m-(2k+1)$ internally vertex-disjoint paths.
 
\medskip\noindent\textbf{Case B: 
$u,v\in R\sigma$ are distinct reflections.} 
This is analogous to Case A. One only needs to swap the roles of rotations and reflections.

\medskip\noindent\textbf{Case C: $u$ is a rotation and $v$ is a reflection (or vice versa).} \\
Without loss of generality, assume $u=\rho^i$ and $v=\rho^j\sigma$ for some $i$ and $j$. 

\begin{itemize}
  \item[] \textsc{Subcase C(A):} Suppose $u$ and $v$ are adjacent. Consider the edge $uv$ as a path.   
  Since $v$ has $m-k$ neighbors in $R$, there are $m-(k+1)$ neighbors of $v$ in $R\backslash \{u\}$.
  For each such rotation $\rho^{r}\in R\setminus\{u\}$ (that is a neighbor of $v$) take the path $P_{\rho^{r}}=\{u\rho^{r},\rho^{r}v\}$.
  Similarly, there are $m-(k+1)$ different reflections distinct from $v$ that are neighbors of $u$. For each such reflection $s$, take the path $P_{s}=\{us, sv\}$. 
  Altogether, these give $1+2(m-(k+1))=2m-(2k+1)$ internally vertex disjoint paths from $u$ to $v$.
  \vspace{2mm}

  \item[] \textsc{Subcase C(B):} If $u$ is not a neighbor of $v$, then similar to \textsc{Subcase C(A)}, we can obtain at least $2m-(2k+1)$ internally vertex disjoint paths from $u$ to $v$.
\end{itemize}

\medskip
By Menger's theorem, $\kappa(X)\ge 2m-(2k+1)=a-(2k+1)> a-b$ since $b=2r-1>2k+1$.
    
    \item Since $X$ is $d$-regular and $a\geq r+8$, we have $\Delta(X)=2m-(k+1)=a-(k+1)\geq r+2$.
\end{enumerate}
\end{proof}

\subsection{circulant Cayley graph}

\begin{prop}\label{Proposition 4.5} 
{\em Let $t=\max\Big\{\Big\lceil\frac{a}{r-1}\Big\rceil-1,\ \Big\lceil\frac{a-b}{2}\Big\rceil,\ \Big\lceil\frac{r+2}{2}\Big\rceil\Big\}$ such that $2t\leq b$.
Let $X=C_a^{\,t}$ be the $t^{th}$-power of the $a$-cycle $C_{a}$. Then $X$ satisfies $(i)-(iii)$ of Theorem \ref{Theorem 4.1}.}
\end{prop}

\begin{proof}
The following shows that $1\le t\le\lfloor \frac{a}{2}\rfloor$ since $a=n-2r+1\geq (3r+3)-2r+1=r+4$.
\begin{enumerate}
    \item We show
$\Big\lceil\frac a{r-1}\Big\rceil-1 \le \Big\lfloor\frac a2\Big\rfloor$ i.e., 
$\Big\lceil\frac a{r-1}\Big\rceil \le \Big\lfloor\frac a2\Big\rfloor +1$.
Since for any real number $x$ we have $\lceil x\rceil \le x+1$, it suffices to verify

\begin{center}
$\frac a{r-1} + 1 \le \Big\lfloor\frac a2\Big\rfloor +1$.   
\end{center}

Since $\lfloor \frac{a}{2}\rfloor \ge \frac{a}{2} - \frac{1}{2}$, it is enough to verify
$
\frac a{r-1} \le \frac{a}{2} - \frac{1}{2}$ i.e., 
$(r-3)a \ge r-1$.

But $r\ge7$ and $a\ge r+4$ give
$(r-3)a \ge (r-3)(r+4) = r^2 + r - 12 \ge r-1
$. 

\item  We show
$\Big\lceil\frac{a-b}{2}\Big\rceil \le \Big\lfloor\frac a2\Big\rfloor.$
Since $b\ge1$, we have
$\frac{a-b}{2}\leq \frac{a}{2} - \frac{1}{2}$. For any natural number $x$, we have
$\lceil \frac{x}{2}-\frac{1}{2}\rceil = \lfloor \frac{x}{2}\rfloor$. Hence,
\begin{center}
$\lceil\frac{a-b}{2}\rceil \le \lceil \frac a2 - \frac12\rceil
=\lfloor\frac a2\rfloor$.    
\end{center}

\item  We show
$
\Big\lceil\frac{r+2}{2}\Big\rceil \le \Big\lfloor\frac a2\Big\rfloor.
$
Since $a\ge r+4$, we have $\lfloor \frac{a}{2}\rfloor \ge \lfloor \frac{r+4}{2}\rfloor$.

\begin{itemize}
  \item If $r=2k$ is even, then $\lceil\frac{r+2}{2}\rceil = k+1$ and
  $\lfloor\frac{r+4}{2}\rfloor = k+2$, so $k+1 < k+2 \le \lfloor \frac{a}{2}\rfloor$.
  \item If $r=2k+1$ is odd, then $\lceil\frac{r+2}{2}\rceil =\lfloor\frac{r+4}{2}\rfloor= k+2\leq \lfloor \frac{a}{2}\rfloor$.
\end{itemize}
\end{enumerate}


If $t=\lfloor \frac{a}{2}\rfloor$, then $X=C_a^{\,t}\cong K_a$.
In view of Example \ref{Example 4.2}, it is enough to assume that $t<\lfloor a/2\rfloor$. 

\begin{enumerate}[(i)]
    \item 
By Lemma \ref{Lemma 2.5}, $\alpha(C_a^{\,t})=\lfloor \frac{a}{t+1}\rfloor$. Since
$t\ge\lceil \frac{a}{r-1}\rceil-1$, we get $\alpha(X)\le r-1$.
   \item Since $t<\lfloor \frac{a}{2}\rfloor$ and $C_a^{\,t}$ is the circulant Cayley graph $\mathrm{Cay}(\mathbb Z_a,\{\pm1,\dots,\pm t\})$, we can see that $C_a^{\,t}$ is a $2t$-regular graph. Thus, $\Delta(X)=2t\ge r+2$ since
$t\ge\lceil(r+2)/2\rceil$. 
\end{enumerate}
 
\begin{claim}\label{Claim 4.6}
    {\em $\kappa(X)\geq a-b$.}
\end{claim}

\begin{proof}
By Menger's theorem, it is enough to show that for an arbitrary pair of vertices $\{u,v\}$, there exists at least $a-b$ internally vertex-disjoint $u$--$v$ paths. 
Let the vertices of $C_a$ be labeled $v_0, v_1, \dots, v_{a-1}$ in clockwise direction. 
Without loss of generality, let $u = v_0$ and $v = v_d$.

\textbf{Case(A): $u$ and $v$ are non-adjacent.}
Then $t < d < a-t$.
We construct $t$ paths moving in the clockwise direction and $t$ paths moving in the anti-clockwise direction.

\textbf{Clockwise Paths ($P_k^+$):}
For each $k \in \{1, 2, \dots, t\}$, let 
\begin{center}
$P_k^+=\{v_{0}v_{k}, v_{k}v_{k+t},v_{k+t}v_{k+2t},...,  v_{k+\lfloor \frac{d-k}{t} \rfloor t} v_d\}$ 
\end{center}

be a path from $v_0$ to $v_d$ that starts with the edge $v_0 v_k$ and the subsequent edges are
$v_{k+(r-1)t}v_{k+rt}$ for $1\leq r\leq {\lfloor \frac{d-k}{t} \rfloor }$.
The edge $v_{k+\lfloor \frac{d-k}{t} \rfloor t} v_{d}$ is valid since $v_{{k+\lfloor \frac{d-k}{t}} \rfloor t}$ is at distance at most $t$ from $v_d$.

\begin{figure}[h]
\centering

\begin{tikzpicture}[scale=1.3]
    % PARAMETERS
    \def\a{12} % total vertices
    \def\t{2}  % hop size
    \def\d{6}  % destination vertex index
%%%%%%%%%%%%%%%%%%%%%%%%%%%%%%%%%%%%%%
% Place vertices around the circle
    \foreach \i in {0,...,\numexpr\a-1} {
        \ifnum\i=0
            \node[circle,draw,inner sep=1pt,label=above:$v_0$] (v0) at ({90-360*\i/\a}:1) {};
        \else
            \node[circle,draw,inner sep=1pt] (v\i) at ({90-360*\i/\a}:1) {};
        \fi
    }
%%%%%%%%%%%%%%%%%%%%%%%%%%%%%%%%%%%%
% Draw the cycle edges with some colored edges
\foreach \i in {0,...,\numexpr\a-1} {
    \pgfmathtruncatemacro{\next}{mod(\i+1,\a)}

    % Conditional coloring
    \ifnum\i=6
        \draw[black] (v\i) -- (v\next);  % 7th edge in blue
    \else\ifnum\i=7
        \draw[very thick,black] (v\i) -- (v\next); % 8th edge in orange
    \else
        \draw[gray!40] (v\i) -- (v\next); % all others gray
    \fi\fi
     \ifnum\i=7
        \node at ($(v\i)+(0,-0.2)$) {$v_{d}$};
    \fi
}
%%%%%%%%%%%%%%%%%%%%%%%%%%%%%%%%%%%%%%%
% ==== CLOCKWISE PATH (blue) =================
    \pgfmathtruncatemacro{\nsteps}{int(\d/\t)}
\foreach \j in {1,...,\nsteps} {
    \pgfmathtruncatemacro{\prev}{mod((\j-1)*\t,\a)}
    \pgfmathtruncatemacro{\idx}{mod(\j*\t,\a)}
    
     % Use dotted line for the 3rd edge
    \ifnum\j=3
        \draw[thick,black,dotted] (v\prev) -- (v\idx);
    \else
        \draw[black] (v\prev) -- (v\idx);
    \fi
    
    \ifnum\j=1
        \node at ($(v\idx)+(0.25,0)$) {$v_{k}$};
    \fi
    % Label the second vertex (v_{k+t})
    \ifnum\j=2
        \node at ($(v\idx)+(0.37,0)$) {$v_{k+t}$};
    \fi
    
    \ifnum\j=3
        \node at ($(v\idx)+(0.45,-0.3)$) {$v_{k+\lfloor \frac{d-k}{t}\rfloor t}$};
    \fi
}
%%%%%%%%%%%%%%%%%%%%%%%%%%%%%%%%%%%
% ==== ANTICLOCKWISE PATH (orange) =================
\pgfmathtruncatemacro{\current}{0} % start at v0
\foreach \i in {1,2} { % only first two edges
    \pgfmathtruncatemacro{\next}{mod(\current - \t + \a,\a)} % always positive

% Use dotted line for the 3rd edge
    \ifnum\i=2
        \draw[very thick,black,dotted] (v\current) -- (v\next);
    \else
        \draw[very thick,black] (v\current) -- (v\next);
    \fi

    \xdef\current{\next} % update current

    \ifnum\i=1
        \node at ($(v\next)+(-0.4,0)$) {$v_{a-k}$};
    \fi
    
    \ifnum\i=2
       \node at ($(v\next)+(-1,0)$) {$v_{a-k-\lfloor \frac{a-d-k}{t} \rfloor t}$};
    \fi
}

\end{tikzpicture}
\caption{\em A clockwise path $P_{k}^{+}$ and an anticlockwise path $P_{k}^{-}$.}
\label{Figure 3}
\end{figure}

\textbf{Anti-Clockwise Paths ($P_k^-$):}
For each $k \in \{1, 2, \dots, t\}$, we similarly construct a path 
\begin{center}
$P_k^-=\{ v_0 v_{a-k}, v_{a-k}v_{a-k-t}, v_{a-k-t}v_{a-k-2t},..., v_{a-k-\lfloor \frac{a-d-k}{t} \rfloor t} v_d\}$     
\end{center}

by moving in the anti-clockwise direction (see Fig. \ref{Figure 3}).

Define $\mathcal{Q}=\{P_{k}^{+}:k\in\{1,...,t\}\}\cup \{P_{k}^{-}:k\in\{1,...,t\}\}$.

\begin{subclaim}\label{Subclaim 4.7}
    {\em The set $\mathcal{Q}$ is a collection of internally vertex disjoint $u-v$ paths.}
\end{subclaim}

\begin{proof}
Consider two distinct clockwise paths $P_k^+$ and $P_j^+$ with $1 \le k < j \le t$. The internal vertices of $P_k^+$ and $P_{j}^{+}$ have indices of the form $k + b_{1}t$ and $j + b_{2}t$ respectively, for some integers $b_{1}, b_{2} > 0$. If they share a vertex, then for some $b_{1}, b_{2}>0$ we have
\[ 
j+b_{2}t \equiv k+b_{1}t \pmod a \text{ i.e., } j-k \equiv (b_{1}-b_{2})t \pmod a. 
\]
Since $1 \le k < j \le t$, we have 
$1\leq j-k\leq t-1$. Also, $a > 2t$. Since the only multiple of $t$ in the range $[-(t-1), t-1]$ is $0$, we obtain $k-j=0$, which contradicts $k \ne j$. Thus, all clockwise paths are internally vertex-disjoint among each other. Similarly, all anti-clockwise paths are also internally vertex-disjoint among each other. 
Furthermore, $P_k^+$ and $P_j^-$ are internally disjoint for any $1 \le k, j \le t$ since the indices of the internal vertices of $P^{+}_{k}$ can be at most $d-1$ where as the indices of the internal vertices of $P^{-}_{j}$ can be at least $d+1$. 

\textbf{Case(B): $u$ and $v$ are adjacent.} Then $\mathcal{R}=\{P_{k}^{-}:k\in \{1,...,t\}\}\cup \{v_{0}v_{d}\}\cup \{\{v_{0}v_{k},v_{k}v_{d}\}:k\in \{1,...,t\}\backslash \{d\}\}$ is a collection of $2t$ internally vertex disjoint $u-v$ path.
\end{proof}

Thus, there are at least $2t$ internally vertex-disjoint paths between the arbitrary chosen pair $\{u,v\}$ of vertices. By Menger's theorem, we have $\kappa(X)\ge 2t\geq a-b$ since $t\ge\lceil(a-b)/2\rceil$. 
\end{proof}

\end{proof}
%%%%%%%%%%%%%%%%%%%%%%%%%%%%%%
\textbf{Funding} The author was supported by the EK\"{O}P-24-4-II-ELTE-996 University Excellence scholarship program of the Ministry for Culture and Innovation from the source of the National Research, Development and Innovation fund. 
%%%%%%%%%%%


\begin{thebibliography}{99}

\bibitem{AS2016}
 S. Alikhani and S. Soltani,
\newblock {The distinguishing chromatic number of bipartite graphs of girth at least six},
\newblock {\em Algebr. struct. appl.}
\newblock \textbf{3}
\newblock (2016),
\newblock 81-87.

\bibitem{BPS2019}
N. Balachandran, S. Padinhatteer, and P. Spiga,
\newblock {Vertex transitive graphs $G$ with $\chi_{D}(G) >\chi(G)$ and small automorphism group},
\newblock {\em Ars Math. Contemp.}
\newblock \textbf{17}
\newblock (2019),
\newblock 311–318. 

\bibitem{BGMa2025} Banerjee, A., Moln\'{a}r, Z.,  Gopaulsingh, A.:
\newblock {Distinguishing chromatic number of middle and subdivision graphs},
\newblock {Transactions on Combinatorics},
\newblock {to appear},
\newblock doi: 
\url{
https://doi.org/10.48550/arXiv.2411.07000} 

\bibitem{BMG2025} Banerjee, A., Moln\'{a}r, Z.,  Gopaulsingh, A.:
\newblock {Upper Bounds for the List-Distinguishing Chromatic Number},
\newblock {Graphs Combin}
\newblock \textbf{41}
\newblock (2025).
\newblock doi: 
\url{https://doi.org/10.1007/s00373-025-02920-x} 

\bibitem{BK1977}
O.V. Borodin and A.V. Kostochka,
\newblock { On an upper bound of a graph’s chromatic number, depending on the graph’s degree and density},
\newblock {\em J. Combin Theory B}
\newblock \textbf{23}
\newblock (1977),
\newblock 247–250.

\bibitem{BKPS2025}
C. Brause, R. Kalinowski, M. Pil\'{s}niak, and I. Schiermeyer,
\newblock {On the distinguishing chromatic number in hereditary graph classes},
\newblock {\href{
https://doi.org/10.48550/arXiv.2505.17193}{arXiv: 2505.17193[math.CO]}}


\bibitem{Cat1978}
P. A. Catlin,
\newblock {A bound on the chromatic number of a graph},
\newblock {\em Discrete Math.}
\newblock \textbf{22}
\newblock (1978),
\newblock 81–83.

\bibitem{Cat1978a}
P. A. Catlin,
\newblock {Another bound on the chromatic number of a graph},
\newblock {\em Discrete Math.}
\newblock \textbf{24}
\newblock (1978),
\newblock 1-6.

\bibitem{CT2006}
K. L. Collins and A. N. Trenk, 
\newblock {The Distinguishing Chromatic Number},
\newblock {\em Electron. J. of Combin.}
\newblock \textbf{13}
\newblock (2006).


\bibitem{Cra2018}
D. Cranston,
\newblock {Proper Distinguishing Colorings with Few Colors for Graphs with Girth at Least 5},
\newblock {\em Electron. J. Combin.}
\newblock \textbf{25}
\newblock (2018),
\newblock $\#$P3.5.

\bibitem{FGHSW2013}
M. Ferrara, E. Gethner, S. G. Hartke, D. Stolee, and P. S. Wenger, 
\newblock {List distinguishing parameters of trees},
\newblock {\em Discrete Appl. Math.}
\newblock \textbf{161}
\newblock (2013),
\newblock 864-869.

\bibitem{FNS2011}
G. Fijav\v{z}, S. Negami, and T. Sano,
\newblock {3-connected planar graphs are 5-distinguishing colorable with two exceptions},
\newblock {\em Ars Math. Contemp.}
\newblock \textbf{4}
\newblock (2011),
\newblock 165–175.

\bibitem{Hae2021}
W. H. Haemers,
\newblock {Hoffman's ratio bound},
\newblock {\em Linear Algebra Appl.}
\newblock \textbf{617}
\newblock (2021),
\newblock 215-219.


\bibitem{IAT2025}
R. Ichishima1, F.A. Muntaner-Batle, Y. Takahashi,
\newblock {On the strength and independence number of powers of
paths and cycles},
\newblock {\em Commun. comb. optim.}
\newblock \textbf{10}
\newblock (2025),
\newblock 717-728.
\newblock{\url{https://doi.org/10.22049/cco.2024.29087.1839}}

\bibitem{LS2009}
C. Laflamme and K. Seyffarth, 
\newblock {Distinguishing chromatic numbers of bipartite graphs},
\newblock {\em Electron. J. Combin.}
\newblock \textbf{16}
\newblock (2009),
\newblock $\#$R76.

\bibitem{Law1978}
J. Lawrence,
\newblock {Covering the vertex set of a graph with subgraphs of smaller degree},
\newblock {\em Discrete
Math.}
\newblock \textbf{21}
\newblock (1978),
\newblock 61–68.

\bibitem{Lov1966}
L. Lov\'{a}sz,
\newblock {On decomposition of graphs},
\newblock {\em Studia Math. Hung.}
\newblock \textbf{1}
\newblock (1966),
\newblock 237–238.

\bibitem{Wat1970}
M. E. Watkins,
\newblock {Connectivity of transitive graphs},
\newblock {\em J. Comb. Theory.}
\newblock \textbf{8}
\newblock (1970),
\newblock 23-29.

\end{thebibliography}
\end{document}
